\section{Exhaustive verification of the finite signing constants}
\label{app:finite-signings}
This appendix gives the computational proof used in
Proposition~\ref{prop:finite-signings}. Label the vertices $0,\ldots,n-1$.
For vertex signs $t_i$, switching replaces $a_{ij}$ by
$a'_{ij}=a_{ij}t_it_j$. Since $Q_{a'}(s)=Q_a(t\cdot s)$, switching
preserves both extrema of $Q$, and hence $R_V$.
Taking $t_0=1$ and $t_j=a_{0j}$ makes every coefficient incident to
vertex zero positive. This normalized representative is unique in its
switching class: a switching that keeps those coefficients positive has
$t_j=t_0$ for all $j$, so it changes no edge. It therefore suffices to
enumerate the $2^{\binom{n-1}{2}}$ choices of negative edges among
vertices $1,\ldots,n-1$: $2,8,64,1024,32768$ representatives for
$n=3,\ldots,7$.

Encode a cut by the subset on side one, with vertex zero always on
side zero. Complementation changes no cut, so these $2^{n-1}$ subsets
include every cut. If a cut has $c$ crossing edges, of which $h$ are
negative, its weight is $c-2h$. The following complete Python 3.10+
program computes the first row of Proposition~\ref{prop:finite-signings}
in exact integer arithmetic, without floating point or a solver. For each
cut, \texttt{mask} records which free edges (those not incident to vertex
zero) cross it, and \texttt{bit\_count()} counts the set bits of an
integer.

\begin{small}
\begin{verbatim}
from itertools import combinations
def min_range(n):
    edges = list(combinations(range(n), 2))
    free = [(i, j) for i, j in edges if i != 0]
    cuts = []
    for side in range(1 << (n - 1)):
        side <<= 1  # Vertex zero stays on side zero.
        cross = lambda i, j: ((side >> i) ^ (side >> j)) & 1
        size = sum(cross(i, j) for i, j in edges)
        mask = sum(cross(i, j) << k
                   for k, (i, j) in enumerate(free))
        cuts.append((size, mask))
    best = len(edges) + 1
    for negative in range(1 << len(free)):
        weights = [size - 2 * (negative & mask).bit_count()
                   for size, mask in cuts]
        best = min(best, max(weights) - min(weights))
    return best
result = [min_range(n) for n in range(2, 8)]
assert result == [1, 2, 4, 4, 5, 8]
print(result)
\end{verbatim}
\end{small}

Every signed cut weight lies between $-\binom n2$ and $\binom n2$,
and the range of cut weights is at most $\binom n2$, because two cut
weights differ by a sum of at most that many coefficients in $\{-1,1\}$.
The initial value \texttt{len(edges) + 1} therefore exceeds every
candidate range. The two loops cover all normalized signings and all cuts,
so their minimum is exactly $\min_a R_V(a)$, not merely an upper bound
from the best witness found. This is an exhaustive finite proof, not an
analytic formula for general $n$.

\begin{table}[htbp]
\centering
\begin{tabular}{clrr}
\toprule
$n$ & Negative edges; all other edges have coefficient $+1$
    & $\min Q$ & $\max Q$\\
\midrule
3 & $\varnothing$ & $-1$ & $3$\\
4 & $\varnothing$ & $-2$ & $6$\\
5 & $12,14,23$ & $-4$ & $4$\\
6 & $14,15,23,25,34$ & $-5$ & $5$\\
7 & $12,13,16,23,25,34$ & $-7$ & $9$\\
\bottomrule
\end{tabular}
\caption{Attaining full-center signings, with vertices numbered from zero.
An entry $ij$ denotes the unordered pair $\{i,j\}$.}
\label{tab:finite-witnesses}
\end{table}

The extrema in Table~\ref{tab:finite-witnesses} can be checked by
evaluating $Q$ on the $2^{n-1}$ vertex-sign vectors with $s_0=1$.
In each case $R_V=(\max Q-\min Q)/2$ equals the computed minimum. Extending the listed $K_6$ witness by vertex six and positive new
edges gives $\min Q=-9$, $\max Q=11$ at the full center, hence ratio
$21/10$; its original six-vertex face retains ratio three.

The accompanying Python 3.10+ script
\path{verification/check_complete_signings.py} implements the same
integer enumeration, records a histogram of all candidate ranges, and
checks the displayed witnesses by evaluating $Q$ on every induced face.
It writes the witnesses, counts, exact rational ratios, and the extended
six-vertex example to \path{verification/complete_signings.json}.
